% ABSTRACT -- draft
The Vera C.\ Rubin Observatory and \textit{Euclid} survey overlapping sky at
complementary wavelengths and resolutions, but their pixel-level data are
rarely combined before catalog construction. We present \jaisp{}, a joint
Rubin--\textit{Euclid} foundation model that ingests six Rubin optical bands
($ugrizy$) together with \textit{Euclid} VIS and NISP ($Y\!J\!H$) imagery at a
common fused resolution, and learns a shared representation by masked
reconstruction. We use the frozen representation to drive two downstream tasks
on $\sim$790 paired tiles in the ECDFS field: source detection with a
center-prediction head, and astrometry via a multi-band joint centroid that
fuses the sharp \textit{Euclid} VIS localization across all instruments.
% --- key results ---
The joint multi-band position reduces per-band centroid scatter relative to a
single-band measurement, with the largest gains at low signal-to-noise
($\sim$2.4--2.5$\times$ over a VIS-only anchor at $\mathrm{SNR}=5$ in a
controlled source-injection test), and is validated against \textit{Gaia} to
$4.4$--$4.8\mas$ at $G<18.5$. We critically examine what the foundation
contributes: through ablations and injection tests we find the cross-instrument
gain is genuine but bounded by the VIS localization floor, and that detection
completeness is suppressed near bright stars by the network's
signal-to-noise input normalization.
% --- OOD ---
Finally we report a first out-of-distribution test, applying the frozen model
without retraining to the EDF-S field: detection and \textit{Euclid} NISP
astrometry transfer cleanly, \todo{state final Rubin optical transfer result
after real-variance rerun}. \todo{one-sentence takeaway on the value of
multi-instrument pixel-level fusion.}
